\section{Space from the Web}

We now start reading the physical world off the engine. The first thing to recover is space itself. In this theory space is not a stage that things sit on. It is a pattern in the web of vibes, and this chapter shows what that means and why it is not a cheat.

\subsection{Distance is counting steps}

We met this in the chapter on graphs, and now it becomes physics. There is no ruler floating over the crystal. The distance between two vibes is just the smallest number of notes you have to cross to get from one to the other. Near means few steps. Far means many. That is the only meaning distance has.

So when you feel that the far wall is ``over there,'' what that comes to, underneath, is that a great many notes separate the vibes making up your eye from the vibes making up the wall. Space is woven from connections. Pull the connections apart and you have not moved through space, you have changed what space is.

\subsection{Why it looks smooth and three-dimensional}

Up close the crystal is a discrete graph, all jagged steps and individual cells. So why does space feel smooth and continuous? For the same reason a photograph looks smooth though it is made of pixels, or a beach looks like a surface though it is made of grains. Step far enough back from a fine-grained thing and the grain disappears into a smooth whole. The crystal's grain is fantastically fine, trillions of trillions of cells, so from our enormous scale it looks like seamless space.

And why three-dimensional? Because, as chapter 10 showed, the crystal can only be a finite-celled hyperbolic crystal in two, three, or four dimensions, and only three holds stable matter. When the simulation grows a crystal and simply measures how many dimensions its step-counting distance behaves like, the answer comes out sharp and stable at the built-in number. Dimension is not painted on. It is a property of how the web is connected, and it can be read straight off the connections.

\subsection{Space is the shape of the relations}

Put the two together and you have the slogan: space is the large-scale shape of the web of vibes, seen from far away. Geometry, the whole subject of distances and shapes and curvature, is a summary of who is connected to whom. The Poincare-disk pictures from Part II were never containers. They were ways of drawing the relations so the shape becomes visible.

This is why moving the drawing changes nothing. Slide the picture, redraw it, and as long as the vibes keep the same neighbors, space is unchanged, because space was the connections, not the drawing. There is no hidden true position underneath. The connections are the whole story.

\subsection{What this buys}

Treating space as emergent rather than fundamental is not just philosophical tidiness. It pays off. Because space is built from a finite, discrete web, there is a smallest length, the size of one cell, below which ``distance'' stops meaning anything. That smallest length is what later tames the infinities that plague ordinary physics, like the infinite densities at the center of a black hole. A world with a smallest length cannot have a true point of infinite anything.

And because the web is hyperbolic and directionless, the emergent space inherits that: no preferred direction, no preferred place, exactly as our universe seems to be. We get the symmetries of space for free, because they were built into the curved web from the start.

\textbf{Takeaway:} space is not a stage but a pattern in the web. Distance is the number of notes between vibes, the smooth look comes from fine grain seen at large scale, and three dimensions is read off the connections. Space is the large-scale shape of who touches whom, which gives a smallest length and the right symmetries for free.
